%% ============================================================
%%  PARTE 3 — Algoritmo Memético: LCA + Solis-Wets
%%  Hasta 1.0 punto
%% ============================================================

\newpage
\section{Parte III: Hibridación Memética — LCA\,+\,Solis-Wets}
\label{sec:parte3}

\subsection{Motivación de la hibridación}

Un \textbf{algoritmo memético} combina una metaheurística poblacional con una
estrategia de búsqueda local, con el objetivo de explotar más eficientemente las
regiones prometedoras del espacio de soluciones. La idea central es que la búsqueda
global del LCA localice cuencas de atracción de buena calidad, y que el refinamiento
local las explore con mayor profundidad de lo que el propio LCA podría hacer por sí solo.

\paragraph{¿Qué es Solis-Wets?}
\textbf{Solis-Wets} es un método clásico de \emph{búsqueda local estocástica} propuesto
por Solis y Wets en 1981 \cite{solis1981minimization}. Opera por muestreo aleatorio
adaptativo: en cada paso genera una perturbación gaussiana $\mathbf{d}\sim\mathcal{N}(0,\sigma)$
y la prueba en torno a la solución actual incorporando un \textbf{sesgo direccional}
$\mathbf{b}$ que \emph{memoriza} las direcciones que han producido mejora (acelerando el
avance por valles), y la dirección opuesta si no la hay. El \textbf{radio $\sigma$ es
adaptativo}: se agranda tras varios éxitos consecutivos y se encoge tras varios fallos,
ajustando automáticamente la escala del refinamiento. Al no usar gradiente, encaja de
forma natural en problemas de caja negra como CEC2017. Es, por tanto, un explotador puro
y barato, idóneo para corregir el cuello de botella diagnosticado en la Parte~II
(explotación local deficiente del LCA base).

La elección de \textbf{Solis-Wets} como método de búsqueda local se justifica por:
\begin{itemize}
  \item Es un método clásico y ampliamente validado para optimización continua.
  \item No requiere información del gradiente, lo que lo hace compatible con las
        funciones de la competición CEC2017 (caja negra).
  \item Su coste por iteración es lineal en $D$, lo que permite aplicarlo de forma
        eficiente sin superar el presupuesto total de evaluaciones.
  \item Ya está disponible en el repositorio \texttt{cec2017real} como ejemplo de referencia.
\end{itemize}

\subsection{Diseño del algoritmo híbrido LCA-SW}

El esquema de hibridación adoptado es el de \textbf{hibridación en cascada periódica}:
cada $k$ generaciones del LCA, se aplica Solis-Wets sobre los $m$ mejores individuos
de la población. Esto permite que la búsqueda global del LCA descubra nuevas regiones,
y la búsqueda local refine las más prometedoras.

\begin{algorithm}[H]
\caption{LCA Memético (LCA-SW)}
\label{alg:lcasw}
\begin{algorithmic}[1]
\Require $N$, $T$, $lb$, $ub$, $k$ (frecuencia de LS), $m$ (individuos refinados), $\delta_0$ (radio inicial SW)
\Ensure $X_{\text{best}}$
\State Inicializar población con ROBL; evaluar fitness \Comment{mismo que LCA base}
\State $X_{\text{best}} \leftarrow$ mejor individuo inicial
\For{$t = 1$ \textbf{to} $T$}
  \State Ejecutar una generación del LCA (Algoritmo~\ref{alg:lca}) \Comment{exploración + explotación}
  \If{$t \mod k = 0$} \Comment{cada $k$ generaciones}
    \State Seleccionar los $m$ mejores individuos de la población
    \For{cada individuo seleccionado $\mathbf{x}_i$}
      \State Aplicar Solis-Wets($\mathbf{x}_i$, $\delta_0$, $\text{maxEvals\_LS}$) \Comment{refinamiento local}
      \State Actualizar $X_{\text{best}}$ si mejora
    \EndFor
  \EndIf
  \State $\zeta \leftarrow t/T$ \Comment{actualizar parámetro adaptativo}
\EndFor
\State \Return $X_{\text{best}}$
\end{algorithmic}
\end{algorithm}

\subsubsection{Gestión del presupuesto de evaluaciones}

Un aspecto crítico en el algoritmo memético es el reparto del presupuesto total de
evaluaciones entre la búsqueda global (LCA) y la búsqueda local (Solis-Wets). En esta
implementación no se reserva un porcentaje fijo a priori: cada búsqueda local tiene un
coste \emph{fijo} de $\text{Eval}_{\text{LS}} = 100$ evaluaciones y se aplica sobre los
$m=3$ mejores individuos cada $k=10$ generaciones. Todas las evaluaciones consumidas por
Solis-Wets se descuentan del mismo contador global \texttt{max\_evals} que la búsqueda
global, de modo que el coste local total es aproximadamente

\begin{equation}
  \text{Coste}_{\text{LS}} \approx \left\lfloor \frac{T}{k} \right\rfloor \cdot m \cdot \text{Eval}_{\text{LS}},
\end{equation}

donde $T$ es el número de generaciones efectivas. Con $k=10$, $m=3$ y
$\text{Eval}_{\text{LS}}=100$, la búsqueda local consume del orden del 15--20\% del
presupuesto, fracción que emerge de los parámetros y no de un reparto impuesto.

\subsection{Implementación del LCA-SW}

La implementaci\'on completa del algoritmo memético ---incluida la rutina Solis-Wets con control adaptativo del radio $\delta$--- se encuentra en el fichero \texttt{testlcasw.cc} del software entregado.

\subsection{Parámetros del LCA-SW}

\begin{table}[h!]
\centering
\caption{Parámetros del algoritmo memético LCA-SW}
\label{tab:params_sw}
\renewcommand{\arraystretch}{1.3}
\begin{tabular}{@{}llll@{}}
\toprule
\textbf{Parámetro} & \textbf{Valor} & \textbf{Descripción} \\
\midrule
$N$ & 30 & Tamaño de la población \\
$k$ & 10 & Frecuencia de búsqueda local (cada $k$ generaciones) \\
$m$ & 3  & Número de individuos refinados por LS \\
$\delta_0$ & 0.4 & Radio inicial de Solis-Wets \\
$\text{Eval}_{\text{LS}}$ & 100 & Evaluaciones por llamada a Solis-Wets \\
\bottomrule
\end{tabular}
\end{table}

\subsection{Resultados del LCA-SW y análisis comparativo}

La Tabla~\ref{tab:comp_sw} compara el error medio del LCA base frente al híbrido memético
LCA-SW sobre una selección representativa de funciones, en las tres dimensiones
(31 ejecuciones por celda).

\begin{table}[h!]
\centering
\caption{Error medio: LCA base vs.\ LCA-SW en $D \in \{10,30,50\}$ (31 ejecuciones)}
\label{tab:comp_sw}
\footnotesize
\renewcommand{\arraystretch}{1.2}
\begin{tabular}{@{}lcccccc@{}}
\toprule
& \multicolumn{2}{c}{\textbf{$D=10$}} & \multicolumn{2}{c}{\textbf{$D=30$}} & \multicolumn{2}{c}{\textbf{$D=50$}} \\
\cmidrule(lr){2-3}\cmidrule(lr){4-5}\cmidrule(lr){6-7}
\textbf{Función} & \textbf{LCA} & \textbf{LCA-SW} & \textbf{LCA} & \textbf{LCA-SW} & \textbf{LCA} & \textbf{LCA-SW} \\
\midrule
F01 (Uni.)   & 2.122e+09 & 4.659e+02 & 1.665e+10 & 7.969e+02 & 4.782e+10 & 5.194e+02 \\
F03 (Uni.)   & 8.545e+03 & 1.713e-01 & 6.203e+04 & 5.563e+04 & 1.309e+05 & 8.227e+04 \\
F04 (Multi.) & 4.711e+01 & 1.831e+00 & 1.365e+03 & 8.481e+01 & 5.528e+03 & 1.504e+02 \\
F11 (Híb.)   & 7.481e+01 & 1.812e+01 & 2.103e+03 & 1.041e+02 & 8.602e+03 & 1.574e+02 \\
F13 (Híb.)   & 3.787e+03 & 2.972e+03 & 2.223e+08 & 1.290e+04 & 4.351e+09 & 1.101e+04 \\
F15 (Híb.)   & 3.260e+02 & 2.814e+02 & 1.351e+06 & 5.686e+03 & 5.119e+08 & 4.457e+03 \\
F25 (Comp.)  & 4.994e+02 & 4.224e+02 & 6.517e+02 & 4.091e+02 & 2.681e+03 & 5.866e+02 \\
F28 (Comp.)  & 6.353e+02 & 4.376e+02 & 1.247e+03 & 4.013e+02 & 3.970e+03 & 5.079e+02 \\
F30 (Comp.)  & 1.426e+05 & 5.845e+04 & 3.574e+07 & 9.849e+05 & 2.989e+08 & 1.985e+07 \\
\bottomrule
\end{tabular}
\end{table}

\begin{table}[h!]
\centering
\caption{Ranking de Friedman (LCA base vs.\ LCA-SW, 30 funciones) y test de Wilcoxon}
\label{tab:sw_stats}
\renewcommand{\arraystretch}{1.3}
\begin{tabular}{@{}lccc@{}}
\toprule
\textbf{Métrica} & \textbf{$D=10$} & \textbf{$D=30$} & \textbf{$D=50$} \\
\midrule
Ranking Friedman LCA base & 1.83 & 2.00 & 1.93 \\
Ranking Friedman LCA-SW   & \textbf{1.17} & \textbf{1.00} & \textbf{1.07} \\
LCA-SW gana (de 30 func.) & 25 & 30 & 28 \\
Wilcoxon SW vs.\ base ($p$) & $0.0005$ & $<10^{-4}$ & $<10^{-4}$ \\
Resultado & SW mejor & SW mejor & SW mejor \\
\bottomrule
\end{tabular}
\end{table}

\subsection{Discusión}

La hibridación memética LCA-SW mejora al LCA base de forma \textbf{estadísticamente
significativa en las tres dimensiones} (Wilcoxon $p < 0.001$), con un ranking de Friedman
prácticamente perfecto a favor de LCA-SW ($\approx 1.0$). En el cómputo global sobre las
30 funciones, LCA-SW gana al LCA base en la inmensa mayoría de casos.

\paragraph{Impacto en funciones unimodales (F1).}
El efecto más dramático se da en F1. El error del LCA base es de $2.1\times 10^9$ ($D=10$),
$1.7\times 10^{10}$ ($D=30$) y $4.8\times 10^{10}$ ($D=50$), reflejo de su incapacidad para
afinar en valles estrechos. Tras incorporar Solis-Wets el error cae a $466$, $797$ y $519$
respectivamente: una mejora de \textbf{7 a 8 órdenes de magnitud}. Esto confirma que la
debilidad del LCA no estaba en la exploración, sino en una explotación local demasiado
grosera, corregida por el paso adaptativo de Solis-Wets.

\paragraph{Escalado con la dimensión.}
A diferencia del LCA base, cuyo error explota al aumentar $D$, el LCA-SW mantiene la mejora
en todas las dimensiones e incluso la acentúa en las funciones híbridas (F11, F13, F15),
donde la búsqueda local recupera varios órdenes de magnitud. La intensificación periódica
sobre los $M=3$ mejores tumores es, por tanto, especialmente valiosa en alta dimensión, donde
el refinamiento manual del LCA resulta más insuficiente.

\paragraph{Ajuste del equilibrio exploración-explotación.}
El éxito de LCA-SW radica en limitar Solis-Wets a un pequeño subconjunto de élites
($M=3$) y solo cada $K=10$ generaciones. Así, la mayor parte de la población sigue
explorando con vuelos de Lévy mientras unas pocas soluciones se refinan en profundidad,
combinando exploración global y explotación local sin agotar el presupuesto ni provocar
convergencia prematura.
